% ===================================================================
%  CUADRO N.º 7 — El Pueblo en invierno. — La Granja en primavera.
%  Traduction 2026 d'apres texto/io, controlee sur texto/fr et
%  texto/en. Consigne : voir texto/es/10-cuadro-01.tex.
% ===================================================================

%%K t07-noto-1 noto
\VUnotes{143.2mm}{%
(*) Para el detalle de una comida, véanse los cuadros 6 y 13.%
}

%%K t07-apar-1 apar
\VUsaut{4.0mm}
\VUcentre{13.2pt}{90}{SEGUNDA SERIE}
\VUsaut{3.0mm}
\VUfilet{16mm}
\VUsaut{5.6mm}
\VUtitre{15.0pt}{60}{CUADRO N.º 7}
\VUsaut{3.0mm}

%%K t07-tit-1 sub
\VUcentre{12.0pt}{0}{\textit{Primera escena.}}
\VUsaut{3.0mm}

%%K t07-tit-2 sub
\VUpk{10.2pt}{40}{El pueblo en invierno.}
\VUsaut{4.0mm}

%%K t07-c1-01-1 p
1. --- Durante las vacaciones de Año Nuevo acompañé a mi padre a
Villanueva, un pueblecito donde tenía algunos asuntos que tratar.

%%K t07-c1-02-1 p
2. --- Tomamos la \VUgras{diligencia} \textsuperscript{(11)} que hace el servicio
entre la ciudad y ese pueblo; pero como hacía mucho frío, aquel viaje en
un vehículo tan poco cómodo no fue nada agradable.

%%K t07-c1-03-1 p
3. --- Por eso sentimos verdadero placer cuando vimos que el
\VUgras{cochero} detenía su carruaje en la plaza, delante de la
\VUgras{posada} \textsuperscript{(2)} del \textit{Oso Blanco}. El \VUgras{posadero}
\textsuperscript{(4)} nos esperaba en el umbral de su casa, fumando tranquilamente
su \VUgras{pipa}.

%%K t07-c1-03-2 p
Nos invitó a entrar y no me hice de rogar para instalarme delante del
fuego de sarmientos que chisporroteaba en el hogar. Entonces me reía del
áspero \VUgras{viento del norte} que hacía chirriar el viejo \VUgras{letrero}
\textsuperscript{(3)} y se llevaba en negros remolinos el \VUgras{humo}
\textsuperscript{(1)} que salía de la chimenea.

%%K t07-c1-04-1 p
4. --- Era la hora de comer. Sin esperar más, hicimos honor a la comida
sencilla pero sabrosa que nos sirvieron (*).

%%K t07-tit-3 sub
\VUpk{10.2pt}{40}{Unas cuantas visitas.}
\VUsaut{3.0mm}

%%K t07-c1-05-1 p
5. --- Después de comer acompañé a mi padre, que es agente de seguros, a
casa de sus clientes. Primero visitamos al \VUgras{herrero} \textsuperscript{(5)},
que vive junto a la posada. Estaba herrando un caballo. Admiré la fuerza
de su ayudante, que sujetaba firmemente contra su delantal de cuero el
\VUgras{casco} del caballo mientras el herrero fijaba la \VUgras{herradura}
\textsuperscript{(6)} a fuertes martillazos.

%%K t07-c1-06-1 p
6. --- Luego fuimos a casa del \VUgras{zapatero} \textsuperscript{(7)} del pueblo, el
tío Jacobo. Aunque tenía por letrero una bota \VUgras{magnífica}, se
contentaba con echar suelas a un par de zapatos viejos y agujereados.

%%K t07-c1-06-2 p
El buen hombre nos dijo riendo: «En la ciudad era “zapatero” y no tenía
clientes. Aquí soy “remendón” y hay trabajo de sobra».

%%K t07-c1-07-1 p
7. --- Nos habló de su vecino el \VUgras{barbero} \textsuperscript{(8)}, cuya
clientela está descontenta. Sus \VUgras{navajas} están siempre melladas y
sus \VUgras{tijeras} no cortan nunca. Pero como es el único barbero del
pueblo, no hay más remedio que dejarse afeitar y cortar el pelo por él.
Además es un hombre avaro y antipático.

%%K t07-c1-08-1 p
8. --- Por suerte, los demás \VUgras{vecinos} no se le parecen. El
\VUgras{carretero} \textsuperscript{(9)}, por ejemplo, es un hombre
excelente, trabajador y servicial. Da gusto llevarle un coche a arreglar.
Su trabajo está siempre bien acabado.

%%K t07-c1-09-1 p
9. --- El tío Jacobo tampoco se quejaba del \VUgras{guarnicionero}
\textsuperscript{(10)}, al que a veces ayuda a coser los \VUgras{arreos} cuando el
trabajo apremia.

%%K t07-c1-09-2 p
Mi padre le preguntó por su familia. «Todos están bien. Mi nieta está en
la \VUgras{escuela} \textsuperscript{(12)}. Su \VUgras{maestra} \textsuperscript{(13)} está
muy contenta con ella; es buena \VUgras{alumna} \textsuperscript{(14)}. No es como
ese granuja de hijo mío; solo piensa en jugar. El \VUgras{maestro}
\textsuperscript{(15)} no consigue enseñarle nada».

%%K t07-tit-4 sub
\VUsaut{3.0mm}
\VUpk{10.2pt}{40}{Charlas del pueblo.}
\VUsaut{3.0mm}

%%K t07-c1-10-1 p
10. --- El buen hombre nos contó luego las novedades del pueblo. Iban a
construir una escuela nueva, porque la que estaba instalada en el
\VUgras{ayuntamiento} se había quedado muy pequeña. Por lo demás era fácil,
porque bastaba con comprar una parte del jardín del \VUgras{molinero}. Eso
le vendría muy bien al pobre hombre. Desde que llegaron los fríos, las
\VUgras{muelas} del \VUgras{molino} \textsuperscript{(16)} ya no podían moler el trigo,
porque la \VUgras{rueda} \textsuperscript{(17)} estaba presa en el \VUgras{hielo}
\textsuperscript{(25)} del \VUgras{arroyo} \textsuperscript{(23)}.

%%K t07-c1-11-1 p
11. --- «Y hablando de construcciones, ¿han visto ustedes nuestra nueva
\VUgras{iglesia} \textsuperscript{(18)}? Está muy pintoresca bajo su manto de nieve.
Los \VUgras{cuervos} \textsuperscript{(19)}, que ya no reconocen su viejo
campanario, se posan tristemente en el gran árbol del \VUgras{cementerio}
\textsuperscript{(20)}».

%%K t07-c1-12-1 p
12. --- La idea del cementerio le recordó al zapatero a las personas que
mi padre había conocido y que habían muerto el año pasado. «¡Cuántas
\VUgras{tumbas} \textsuperscript{(21)}, mi buen señor, se han añadido a las otras
bajo el \VUgras{ciprés} \textsuperscript{(22)}! La muerte se llevó a nuestro viejo
notario. Todo el pueblo asistió a su \VUgras{entierro}».

%%K t07-c1-13-1 p
13. --- «Ahí está su sucesor, patinando en el arroyo. Vino hace un rato a
que le arreglara la correa de su \VUgras{patín} \textsuperscript{(26)}, que se había
roto».

%%K t07-c1-14-1 p
14. --- «Ahí está también, a la orilla del arroyo, el nuevo
\VUgras{guarda rural} \textsuperscript{(27)}. Es un antiguo gendarme y el terror de
los golfillos del pueblo. Huyen en cuanto ven sus \VUgras{polainas}
\textsuperscript{(29)} y su \VUgras{quepis} \textsuperscript{(28)}».

%%K t07-c1-15-1 p
15. --- «¡Ah! Ahí está el tío José, que trae un \VUgras{carro de haces de
leña} \textsuperscript{(30)} para el panadero. --- Es verdad, dijo mi padre,
reconozco su tiro de \VUgras{mulas} \textsuperscript{(31)}. Tengo que hablar con él y
voy a aprovechar la ocasión. Hasta la vista, tío Jacobo».

%%K t07-c1-16-1 p
16. --- Justo cuando salíamos, los niños del pueblo dejaban la escuela y
se lanzaban corriendo al \VUgras{resbaladero} \textsuperscript{(33)}, con riesgo de
caerse y romperse un miembro en la \VUgras{caída} \textsuperscript{(34)}. Uno de
ellos cogió su \VUgras{trineo} \textsuperscript{(32)} en casa del carretero, sentó
dentro a uno de sus compañeros y partió corriendo.

%%K t07-c1-17-1 p
17. --- Otros se precipitaron hacia un \VUgras{montón de nieve}
\textsuperscript{(37)} junto al \VUgras{puente} \textsuperscript{(40)} y empezaron a
bombardear al \VUgras{muñeco de nieve} \textsuperscript{(39)} que habían hecho el día
antes y al que habían puesto un sombrero, una pipa y una cartera de
colegio en bandolera.

%%K t07-c1-18-1 p
18. --- Poco les importaban el viento helado y el frío. La mayoría ni
siquiera tenía \VUgras{bufanda} \textsuperscript{(35)} y, para entrar en calor
después de la batalla de \VUgras{bolas de nieve} \textsuperscript{(38)}, les bastaba
un puñado de \VUgras{castañas} \textsuperscript{(42)} muy calientes compradas a la
tía Jacoba, la \VUgras{vendedora}, que tirita de frío bajo su \VUgras{chal}
\textsuperscript{(43)} demasiado delgado.

%%K t07-tit-5 sub
\VUsaut{3.0mm}
\VUcentre{12.0pt}{0}{\textit{Segunda escena.}}
\VUsaut{3.0mm}

%%K t07-tit-6 sub
\VUpk{10.2pt}{40}{La casa de labor en primavera.}
\VUsaut{4.0mm}

%%K t07-c2-01-1 p
1. --- A pesar de todos los placeres del invierno, con hondo gozo
saludamos la vuelta de la \VUgras{primavera}.

%%K t07-c2-01-2 p
¡Qué cuadro tan alegre es el corral de una granja, al atardecer, en el
mes de mayo!

%%K t07-c2-02-1 p
2. --- En el horizonte, el \VUgras{sol} \textsuperscript{(1)} se pone despacio,
inundando el \VUgras{campo} \textsuperscript{(3)} con sus \VUgras{rayos}
\textsuperscript{(2)} de púrpura. El diligente \VUgras{labrador} \textsuperscript{(6)} se da
prisa en arar su \VUgras{campo} \textsuperscript{(4)}.

%%K t07-c2-02-2 p
Con su \VUgras{arado} \textsuperscript{(5)}, tirado por un vigoroso \VUgras{caballo}
\textsuperscript{(7)}, traza el \VUgras{surco} \textsuperscript{(8)} en el que el
\VUgras{sembrador} \textsuperscript{(9)} echará a puñados la \VUgras{simiente} que dará
las espigas doradas.

%%K t07-c2-03-1 p
3. --- Los \VUgras{segadores} \textsuperscript{(11)}, con sus guadañas, han cortado la
hierba del prado, los henificadores han secado el \VUgras{heno}
\textsuperscript{(10)} volviéndolo con las \VUgras{horcas} \textsuperscript{(12)} y los
rastrillos, y helos aquí de vuelta.

%%K t07-c2-03-2 p
Unos van delante del pesado carro tirado por bueyes; llevan su herramienta
al hombro. Otros, más perezosos, se han instalado cómodamente sobre el
heno oloroso.

%%K t07-c2-04-1 p
4. --- Al pasar saludan amistosamente al \VUgras{jardinero} \textsuperscript{(16)},
ocupado en el \VUgras{huerto de frutales} \textsuperscript{(13)} en podar los
\VUgras{árboles frutales} e injertar los \VUgras{perales} \textsuperscript{(14)}. Los
\VUgras{ciruelos} \textsuperscript{(15)} están en flor y, en la \VUgras{huerta}
\textsuperscript{(17)} bien abonada, las verduras, coles y lechugas ya están
hermosas.

%%K t07-c2-04-2 p
Sobre el murete, un hermoso \VUgras{pavo real} \textsuperscript{(18)} despliega la
cola para que admiren sus magníficas plumas.

%%K t07-tit-7 sub
\VUpk{10.2pt}{40}{El corral de la granja.}
\VUsaut{3.0mm}

%%K t07-c2-05-1 p
5. --- Las puertas de la \VUgras{cuadra} \textsuperscript{(19)} están abiertas y se ven
el pesebre y la \VUgras{cama de paja} \textsuperscript{(23)} de la \VUgras{yegua}
\textsuperscript{(21)} que el \VUgras{mozo de cuadra} \textsuperscript{(20)} acaba de
desenganchar. El \VUgras{potrillo} \textsuperscript{(22)}, tan gracioso con sus patas
largas, sigue a su madre dando brincos.

%%K t07-c2-06-1 p
6. --- Junto al \VUgras{pozo} \textsuperscript{(29)}, las \VUgras{vacas}
\textsuperscript{(25)} van a beber en el \VUgras{pilón} \textsuperscript{(26)} de madera que
les sirve de abrevadero. El \VUgras{ternero} \textsuperscript{(24)} aprovecha para
mamar, y el \VUgras{buey} \textsuperscript{(27)}, oliendo el buen olor del forraje,
muge alegremente y levanta con orgullo la cabeza de poderosos
\VUgras{cuernos} \textsuperscript{(28)}.

%%K t07-c2-07-1 p
7. --- Todos trabajan. La \VUgras{criada} \textsuperscript{(32)} sostiene con la mano
la \VUgras{cuerda} \textsuperscript{(31)} del pozo. Saca agua con un \VUgras{cubo}
\textsuperscript{(30)} para dar de beber al ganado. Junto al \VUgras{granero}
\textsuperscript{(33)}, un hombre rastrilla las briznas de paja, y delante de la
puerta de la \VUgras{casa} \textsuperscript{(34)}, una vieja \VUgras{hilandera}
\textsuperscript{(35)}, con su rueca y su \VUgras{copo}, hila el \VUgras{lino} o el
\VUgras{cáñamo} como en otros tiempos.

%%K t07-tit-8 sub
\VUsaut{3.0mm}
\VUpk{10.2pt}{40}{El granjero.}
\VUsaut{3.0mm}

%%K t07-c2-08-1 p
8. --- El \VUgras{granjero} \textsuperscript{(36)} dirige todo este trabajo y da
órdenes a sus criados. Manda recoger la \VUgras{grada} \textsuperscript{(40)} y el
\VUgras{pico} \textsuperscript{(39)} que se han olvidado junto a la \VUgras{caseta}
\textsuperscript{(38)} del \VUgras{perro} \textsuperscript{(37)}.

%%K t07-c2-09-1 p
9. --- Sus gritos espantan a una \VUgras{paloma} \textsuperscript{(41)}, que huye a
todo vuelo hacia el \VUgras{palomar} \textsuperscript{(42)}, y a las \VUgras{gallinas}
\textsuperscript{(43)}, que se apresuran a volver al \VUgras{gallinero}
\textsuperscript{(44)}.

%%K t07-c2-10-1 p
10. --- Delante del \VUgras{redil} \textsuperscript{(48)}, he aquí al viejo
\VUgras{pastor} \textsuperscript{(45)} que trae de vuelta su \VUgras{rebaño}
\textsuperscript{(49)}. Con su \VUgras{cayado} \textsuperscript{(46)}, que no lo deja nunca,
igual que su \VUgras{blusón} \textsuperscript{(47)} descolorido, lleva al aprisco
\VUgras{carneros} \textsuperscript{(50)}, \VUgras{ovejas} \textsuperscript{(53)},
\VUgras{corderos} \textsuperscript{(54)}, \VUgras{cabras} \textsuperscript{(51)} y
\VUgras{cabritos} \textsuperscript{(52)}. Su fiel \VUgras{perro} \textsuperscript{(55)} Argos
cierra la marcha y anima con sus ladridos a los rezagados.

%%K t07-c2-11-1 p
11. --- Todo este ruido y este movimiento no turban la calma de la buena
\VUgras{vaca lechera} \textsuperscript{(56)}, que hace sonar su \VUgras{esquila}
\textsuperscript{(57)} mientras la ordeñan delante de la puerta del
\VUgras{establo} \textsuperscript{(58)}.

%%K t07-c2-12-1 p
12. --- Parece comprender que tiene que dar su leche para que la
\VUgras{granjera} \textsuperscript{(60)} pueda hacer \VUgras{mantequilla} en la
\VUgras{mantequera} \textsuperscript{(61)} o los excelentes \VUgras{quesos}
\textsuperscript{(64)} que escurren sobre la tabla puesta encima del
\VUgras{cuezo} \textsuperscript{(63)}.

%%K t07-c2-13-1 p
13. --- Toda la \VUgras{lechería} \textsuperscript{(59)} la mantiene muy limpia el
\VUgras{mozo de la granja} \textsuperscript{(65)}, que acaba de lavar los
\VUgras{cántaros de la leche} \textsuperscript{(62)} en el gran cubo que lleva en la
mano.

%%K t07-tit-9 sub
\VUsaut{3.0mm}
\VUpk{10.2pt}{40}{El corral de las aves.}
\VUsaut{3.0mm}

%%K t07-c2-14-1 p
14. --- Con sus gruesos zuecos anda con algo de torpeza, y así hace huir
al \VUgras{pavo} \textsuperscript{(68)}, a las \VUgras{patas} \textsuperscript{(66)}, a los
\VUgras{patos} \textsuperscript{(67)} y a los \VUgras{gansos} \textsuperscript{(69)}, que abren
mucho el \VUgras{pico} \textsuperscript{(70)} y lanzan gritos inquietos.

%%K t07-c2-14-2 p
El \VUgras{gallo} \textsuperscript{(71)}, más belicoso, levanta su \VUgras{cresta}
\textsuperscript{(72)} roja, sacude las plumas de su \VUgras{cola} \textsuperscript{(73)} y
lanza un sonoro «quiquiriquí».

%%K t07-c2-15-1 p
15. --- Una \VUgras{gallina madre} \textsuperscript{(74)} picotea en el
\VUgras{estercolero} \textsuperscript{(81)} con sus \VUgras{pollitos} \textsuperscript{(75)}.
El \VUgras{lechón} \textsuperscript{(78)} corre hacia su madre, la enorme
\VUgras{cerda} \textsuperscript{(77)} que vemos junto a la pocilga.

%%K t07-c2-16-1 p
16. --- En sus jaulas, los \VUgras{conejos} \textsuperscript{(79)} comen con avidez la
\VUgras{hierba} \textsuperscript{(80)} tierna que les trae la hija del granjero.
Juanita quiere mucho a sus conejos y, todos los días, después de ir a
buscar los \VUgras{huevos} \textsuperscript{(76)} frescos al gallinero, viene a
visitar a sus amiguitos.
\VUsaut{6.0mm}
\VUfilet{22mm}
